%% conformal-antenna-diffopt.6 -- revision of v5 against the v5 referee review (27/44),
%% under the operator's rulings on the companion transmon paper (gradient advantage
%% demonstrated here; no "certified"; short Section 9). Adds the forward-accuracy study,
%% curvature, power balance and gain, gradient-free baselines with geode-fem and with
%% Palace in the loop, a realizable feed and the design re-centred on the refined model.
%% See changelog.md, provenance.md and verify_evidence.py. v1-v5 and their critic
%% siblings are unchanged.

\documentclass{anvil-paper}
\usepackage{xurl}
\usepackage{array}

\expandafter\def\expandafter\UrlBreaks\expandafter{\UrlBreaks\do\-\do\_}
\Urlmuskip=0mu plus 1mu
\renewcommand{\topfraction}{0.92}
\renewcommand{\bottomfraction}{0.6}
\renewcommand{\textfraction}{0.07}
\renewcommand{\floatpagefraction}{0.85}
\setcounter{topnumber}{3}
\setcounter{totalnumber}{4}
\renewcommand{\bibfont}{\small}
\setlength{\bibsep}{2pt plus 1pt}

\newcommand{\anvilfig}[3][0.98]{%
  \IfFileExists{figures/#2}%
    {\includegraphics[width=#1\linewidth]{figures/#2}}%
    {\fbox{\parbox{#1\linewidth}{\centering\vspace{5em}%
      \texttt{figures/#2}\\[0.5em]%
      {\small placeholder: render via \texttt{figures/src/#3}}%
      \vspace{5em}}}}}
\newcommand{\PL}{P2/L1}

\title{An Adjoint Gradient Reaches the Band Target with an Order of Magnitude
Fewer Solves than Gradient-Free Search: Copper-Only Design of a 2.4\,GHz
Conformal Patch Antenna}
\author{Robb Walters\\
  {\normalsize Independent Researcher}}
\date{October 2026}

\begin{document}
\maketitle

\begin{abstract}
We design the copper of a probe-fed 2.4\,GHz patch antenna on a standard
$1.6$\,mm FR-4 board bent to a $100$\,mm radius, with 38 shape parameters
and a finite-element solver, geode-fem, whose discrete shape adjoint gives
the gradient of the band objective at one factorization per frequency. On
the coarse optimization model the adjoint reaches a worst-of-band
$|S_{11}|$ of $-9.5$\,dB in $35$ solves. Separable CMA-ES and NEWUOA,
driven by the independent solver Palace, need $385$ to $805$ solves
($11$--$23\times$) and $7.6$ to $15.5$ times the wall-clock on the same host.
Palace's forward solve is about $11\%$ faster than geode-fem's, so the
advantage comes from the gradient. The 38 parameters are not needed to meet
the target: a four-parameter rectangle, given the final frequency correction
in advance, also meets it on the second-order model it was corrected for.
The forward model's error sits at the patch rim, and refining the rim in geode-fem and raising the element order in Palace
meet at the same resonance. After one more frequency space-mapping step,
the design with a $1.6$\,mm probe reaches a worst in-band $|S_{11}|$ of
$-10.46$\,dB on Palace with third-order elements, with a $-10$\,dB band of
$2.3935$--$2.4864$\,GHz ($93.0$\,MHz) that covers the whole $2.400$--$2.4835$\,GHz
ISM band with margins of $6.5$ and $2.9$\,MHz. geode-fem's finest mesh agrees
on the lower edge to $0.02$\,MHz. The lower margin covers the order-1
Richardson estimate of the remaining upward drift ($6.4$\,MHz) but not a
rim-size fit ($13.5$\,MHz), at which $2.400$\,GHz would sit at $-8.94$\,dB.
The mesh-morph volume guard, not the objective, ends the final descent and
so limits the bandwidth.
\end{abstract}

\section{Introduction}
\label{sec:intro}

Gradient-based inverse design is routine in nanophotonics, where
differentiable solvers such as \textsc{ceviche} \citep{hughes2019forward},
the Meep adjoint solver \citep{hammond2022meepadjoint}, SPINS-B
\citep{su_spinsb} and FDTDX \citep{mahlau2026fdtdx} optimize a permittivity
density on a structured grid. Shape derivatives of scattering parameters for
finite-element models with ports are older
\citep{webb2000design,webb2002design,nikolova2004adjoint} and are available
in commercial tools
\citep{ansys_hfss_derivatives,comsol_sensitivity_method}. A gradient is worth
its implementation cost only if it buys something a gradient-free method
does not. For a design with two nearly independent parameters, finite
differences or a scaling law are cheap and a gradient adds little. The case
for a gradient is a design with many coupled parameters, where finite
differences cost one solve per parameter and gradient-free search scales
poorly with dimension.

This paper measures that case on one antenna. A 2.4\,GHz probe-fed patch is
bent around a cylinder of radius $R_b = 100$\,mm. The board is a uniform
$1.6$\,mm FR-4 shell that is not designed; 38 smooth parameters move the
copper outline, the corners and the feed. geode-fem, a differentiable
finite-element solver with a discrete shape adjoint, optimizes the band
objective by L-BFGS. We run the same problem, objective, parameters and
volume guard with two gradient-free optimizers, separable CMA-ES
\citep{hansen2001cmaes,ros2008separable} and NEWUOA
\citep{powell2006newuoa,zhang2023prima}, each with geode-fem and with Palace
\citep{palace} as the forward solver, and count solves and wall-clock to fixed
objective levels on one host.

Three results follow.
\begin{itemize}
  \item \textbf{The gradient makes the search about an order of magnitude
  cheaper} (Section~\ref{sec:baselines}). To reach $-9.5$\,dB on the
  optimization model the adjoint uses $35$ solves; the gradient-free runs use
  $385$ to $805$ ($11$--$23\times$) and $7.6$ to $15.5$ times the wall-clock.
  Palace's forward evaluation is cheaper than geode-fem's, so the difference
  is the gradient, not the solver. Given enough solves, one CMA-ES seed of
  three finds a design as good as the adjoint's, and a four-parameter
  rectangle also meets the target on the second-order model, so the gradient
  buys speed rather than a better optimum, and the 38-parameter basis is not
  what makes the target reachable.
  \item \textbf{The forward model's error sits at the patch rim}
  (Section~\ref{sec:accuracy}). Refining the rim moves the resonance;
  refining the substrate interior does not. Palace on geode-fem's own
  meshes reproduces geode-fem at first and second order, and Palace at third
  order agrees with geode-fem's finest rim-refined mesh to under $1$\,MHz.
  \item \textbf{A design re-centred on the refined model meets the ISM band at
  $-10$\,dB on both fine models} (Section~\ref{sec:results}), with a
  probe of realizable width ($1.6$\,mm) and margins small enough that we state which
  drift estimates they cover.
\end{itemize}
We also report the checks around the design. The bend moves the resonance
by under $0.5$\,MHz, though it changes the match by about 1\,dB; the
absorber and air margin move the band by at most a few megahertz
(Section~\ref{sec:accuracy}); and the radiation efficiency of a closely
related design is $0.62$ with perfectly conducting copper
(Section~\ref{sec:results}).

\paragraph{Evidence tiers.}
We keep three tiers apart. A \emph{bounded} result is a proven two-sided
bound; no result in this paper is of that kind. A \emph{model-dependent}
result is a value computed for a stated discrete model (mesh, element order,
absorber and port model). An \emph{estimated} result is an extrapolation to
the mesh limit, given with its band and never called a converged value.

\section{Related Work}
\label{sec:related}

\paragraph{Finite-element S-parameter sensitivities with ports.}
Design sensitivities of the scattering matrix for tetrahedral edge elements
were derived by \citet{webb2000design}, extended to the frequency response
by \citet{webb2002design} and used in gradient optimization of microwave
devices by \citet{nair2003optimization}. Adjoint and self-adjoint
formulations are reviewed by \citet{nikolova2004adjoint} and
\citet{nikolova2006network}, and continuum shape sensitivity for
high-frequency devices is given by \citet{choi2011generalized}. HFSS computes
S-parameter derivatives with respect to geometric parameters by reusing the
LU factors \citep{ansys_hfss_derivatives,vardapetyan2008sensitivity}, and
COMSOL documents adjoint shape optimization of waveguide devices on a
deformed mesh \citep{comsol_sensitivity_method,jensen2022comsolshape}. For
antennas, \citet{ghassemi2013antenna} evolve a microstrip antenna with an FEM
adjoint sensitivity and a small set of control vertices.

\paragraph{Variable-fidelity antenna optimization and space mapping.}
Space mapping aligns a cheap coarse model with an accurate fine model and
optimizes the coarse one \citep{bandler1994space,bandler2004space,koziel2006framework}.
In antenna design the coarse model is often the same electromagnetic solver
on a coarse discretization, corrected by frequency scaling or by
response-shape correction \citep{koziel2010sprp,koziel2014antenna}, and
adjoint sensitivities have been used inside trust-region space mapping
\citep{koziel2013reliable}. Our correction is the simplest member of this
family: a frequency shift fitted between two discretizations of the same
solver.

\paragraph{Automated shape differentiation.}
Shape derivatives of finite-element models can be generated automatically:
in UFL and Firedrake \citep{ham2019automated,dokken2020,paganini2021fireshape}
and in NGSolve \citep{gangl2021fully}. \citet{wang2011} give a discrete
Maxwell shape adjoint on unstructured meshes in two dimensions, in the time
domain and without ports. geode-fem obtains the element Jacobians by
forward-mode differentiation of its closed-form element kernels inside a
tensor stack \citep{burn}.

\paragraph{Structured-grid design of metallic antennas.}
Density-based topology optimization of conductors on a fixed grid goes back
at least to \citet{erentok2011topology}; \citet{hassan2014topology} and
\citet{hassan2014patch} optimize the patch and ground of coax-fed microstrip
antennas with an FDTD adjoint, and
\citet{nikolova2004sensitivity,nikolova2006scattering} give FDTD shape
sensitivities. A plain Yee grid staircases a curved conductor
\citep{cangellaris1991analysis,haggblad2014accuracy}, but locally conformal
FDTD handles curved PEC \citep{dey1997locally,mittra2000conformal},
including uniformly stable second-order schemes
\citep{zagorodnov2003uniformly,benkler2006new}. Subpixel-smoothed projection
gives the Meep density adjoint subpixel-accurate boundaries
\citep{hammond2025unifying,romano2026ssp2}, though the material-grid adjoint
excludes metals \citep{meepadjointdocs}; Tidy3D offers autograd shape
gradients and conformal PEC averaging \citep{tidy3dautograd}; and
\citet{hooten2025} and \citet{liu2024imagerep} map shape parameters to a
fixed grid. Boundary-integral methods carry recent PEC shape-sensitivity
work \citep{takahashi2025tdbie,arens2026tubular,balouchev2024balun}.

\paragraph{Conformal and cylindrical patch antennas.}
The cylindrical-rectangular patch has a cavity-model analysis
\citep{luk1989cylindrical}, and nonplanar microstrip antennas and conformal
arrays are treated in the monographs of \citet{wong1999nonplanar} and
\citet{josefsson2006conformal}. We use a full-wave model and do not compare
with these analyses; the curvature here is mild (Section~\ref{sec:accuracy}).

\paragraph{Gradient-free optimizers.}
CMA-ES adapts a full covariance matrix \citep{hansen2001cmaes}; its
separable variant restricts it to the diagonal, which makes each update
linear in the dimension \citep{ros2008separable}. NEWUOA builds a quadratic
interpolation model inside a trust region \citep{powell2006newuoa}; we use
the PRIMA implementation \citep{zhang2023prima}.

\paragraph{Position of this work.}
The closest prior work is the FEM S-parameter sensitivity line and the
variable-fidelity antenna literature, which already reuse factorizations,
differentiate with respect to geometry and correct coarse models. What this
paper adds is a measurement: on a 38-parameter open-radiator problem with a
moving feed, the solve and wall-clock cost of the adjoint against two
gradient-free optimizers driven by two independent solvers, and a design
checked on two fine models, one of them from an independent solver. We do not claim a first EM shape
adjoint, and we make no performance comparison with structured-grid methods.

\section{Method}
\label{sec:method}

\begin{figure}[t]
  \centering
  \anvilfig{fig1_geometry.pdf}{fig1_geometry.py}
  \caption{The fixture and the final copper design. (a) Section at $y = 0$,
  to scale: the FR-4 board (green) on the $R_b = 100$\,mm cylinder with the
  patch on top, inside the air box and the $25$\,mm box UPML (hatched).
  (b) Section through the feed at $s = 0$, vertical axis stretched four
  times: patch, $1.6$\,mm FR-4, ground, and the $1.6$\,mm probe, final
  (filled) and initial (dashed). (c) Top view in the chart $(s, y)$: the
  initial $38\times29$\,mm patch (dashed), the final copper (filled) and the
  probe positions (bars). The board is the same in both. Drawn from
  \texttt{baseline\_2g4\_L1.toml} and
  \texttt{recentre/geometry/rbG\_L1\_copper\_38\_sm57\_w16.json}.}
  \label{fig:geometry}
\end{figure}

\subsection{The open-radiator forward model}
\label{sec:pencil}

The forward problem is the driven frequency-domain curl-curl equation,
discretized with N\'ed\'elec edge elements of order $p = 1$ or $p = 2$ on an
unstructured tetrahedral mesh. Lengths are in millimetres and the solver
works with $\omega = k_0 = 2\pi f/c$ in rad/mm. The assembled system is the
complex-symmetric pencil
\begin{equation}
  A(\omega)\,\mathbf{x}
  = \Big[\,K(\nu) - \omega^2 M(\varepsilon) + \tfrac{j\omega}{Z_s}\,S_p\,\Big]\,\mathbf{x}
  = \mathbf{b},
  \label{eq:pencil}
\end{equation}
with curl-curl stiffness $K$, mass matrix $M$ and lumped-port surface term
$S_p$. The radiating boundary is a box UPML of thickness $t = 25$\,mm:
per-element tensor weights $\Lambda$ enter $K$ (as $\nu = \Lambda^{-1}$) and
$M$ (as $\varepsilon = \varepsilon_r \Lambda$), with the stretch
$\varsigma = 1 - j\sigma_0 (d/t)^2/\omega$ per axis at depth $d$ and
$\sigma_0 = 1$. The FR-4 has $\varepsilon_r = 4.4$ and $\tan\delta = 0.02$,
and the conductors are perfect electric conductors (PEC). The port is a
rectangle of width $w_p$ along $y$ and length $l_p = 1.6$\,mm across the
substrate, with sheet impedance $Z_s = R_0\,w_p/l_p$ and $R_0 = 50\,\Omega$.
The reflection coefficient follows from the port voltage $V$:
\begin{equation}
  S_{11} = \frac{Z - R_0}{Z + R_0},\qquad Z = \frac{V}{I},\qquad
  I = \frac{2V_{\mathrm{inc}} - V}{R_0}.
  \label{eq:s11}
\end{equation}
Every evaluation checks passivity, $|S_{11}| \le 1$.

\subsection{Shape adjoint with one factorization per frequency}
\label{sec:adjoint}

The design moves mesh node coordinates $X$, so the gradient needs
$\partial K/\partial X$, $\partial M/\partial X$, $\partial S_p/\partial X$
and $\partial \mathbf{b}/\partial X$. geode-fem evaluates its closed-form
element kernels in first-order dual arithmetic: seeding the tangent of one
node coordinate returns the exact partial derivatives of every local matrix
entry. The UPML weights are held fixed under node motion. Because $A(\omega)$
is complex symmetric, the adjoint system uses the same operator, and one LU
factorization per frequency serves the forward and the adjoint solve. The
nodal gradient is $\lambda^{\top}(\partial A/\partial X)\,\mathbf{x}$ plus the
explicit dependence of the port drive and of the port-voltage functional on
the geometry. When a parameter moves the feed, the port faces move and the
port is rebuilt on every moved mesh. A unit test of the moving-feed gradient
on a curved mesh, with a source that varies across the port so that the
explicit term matters, matches central finite differences to below
$10^{-8}$; without the explicit term it misses by a relative error of about 3.

\subsection{Copper design freedom on a fixed curved board}
\label{sec:freedom}

The board is not a design variable. Each substrate node has chart
coordinates $(s, y)$, with $s = R_b\,\mathrm{atan2}(x, z - z_{\mathrm{axis}})$
the mid-plane arc length and $y$ the bend axis. A design applies a chart
displacement $(\Delta s, \Delta y)$ and moves the node to angle
$\phi_0 + \Delta s/R_b$ and height $y_0 + \Delta y$ at its own radius, so
the FR-4 stays a uniform $1.6$\,mm shell. Every run records the largest
change of any substrate node's radius ($\le 1.4\times10^{-14}$\,mm). Air
nodes follow a harmonic extension of the motion; the UPML and the outer
boundary are fixed. The design gradient chains the nodal adjoint gradient
through the exact Jacobian of this map.

Table~\ref{tab:dofs} lists the 38 parameters. They are smooth functions of
the chart, so the same coefficient vector defines the same copper on every
mesh of the family, which lets a design be carried between discretizations.
The basis cannot open a hole in the patch; slots change the conductor's
topology and are outside it.

\begin{table}[t]
  \centering
  \caption{The 38 copper parameters, per unit coefficient. The board
  thickness has none.}
  \label{tab:dofs}
  \small
  \begin{tabular}{@{}l c >{\raggedright\arraybackslash}p{0.66\linewidth}@{}}
    \toprule
    Group & Count & What moves \\
    \midrule
    Outline & $32$ & $8$ clamped cubic B-splines per patch edge ($\pm y$, $\pm s$); each moves its edge outward in the surface by $1$\,mm, ramped in from the feed band and decaying over the $12$\,mm substrate margin \\
    Corners & $4$ & one bump per corner, moving it outward along the bisector by $1$\,mm (negative values chamfer), support $1 - \mathrm{smoothstep}(d/6\,\mathrm{mm})$ \\
    Feed & $2$ & the probe's position along $y$ (its inset from the $+y$ edge) and its width $w_p$, each $1$\,mm \\
    \bottomrule
  \end{tabular}
\end{table}

\subsection{Objective, optimizer and guards}
\label{sec:objective}

The objective is a soft maximum of $|S_{11}|^2$ over $n_f = 5$ band
frequencies,
\begin{equation}
  G(\mathbf{c}) = \frac{1}{\beta}\ln \sum_{f} \exp\!\big(\beta\,|S_{11}(f;\mathbf{c})|^2\big),
  \qquad \beta = 25,
  \label{eq:objective}
\end{equation}
over the coefficient vector $\mathbf{c}$. $G$ exceeds the largest
$|S_{11}|^2$ by at most $\ln 5/\beta = 0.064$, and every level reported in
this paper is the worst of the band points, not $G$. L-BFGS (memory 8) with
Armijo backtracking and a cap of 1 coefficient unit per step minimizes $G$.
A per-element guard rejects any trial step that would shrink a
tetrahedron's signed volume below $0.25$ of its base value. Every descent reported
here ends on the guard with the gradient still nonzero (one intermediate
re-based run of the space-mapping chain ended on a plateau), so the designs
are guard-limited rather than stationary.

Before and after each run the pipeline gradient is compared with a central
finite difference along the steepest-descent direction and along random
directions. The worst relative error over the runs reported here is
$1.5\times10^{-5}$ (tolerance $5\times10^{-3}$), and a per-group check on the
two coarsest meshes passes to $1.9\times10^{-7}$. The adjoint is implemented
for $p = 1$ only, so optimization runs at $p = 1$ and finer models enter
through verification.

\section{Fixture and Forward-Model Accuracy}
\label{sec:accuracy}

\paragraph{Antenna.}
The initial antenna is a $38\times29$\,mm PEC patch (width along $s$,
resonant length along $y$) on $1.6$\,mm FR-4 with a $12$\,mm substrate and
ground margin, fed by a $1.6$\,mm probe, modelled as a uniform lumped port,
inset $7.0$\,mm from the $+y$ edge (Figure~\ref{fig:geometry}). The bend axis
is parallel to $y$, so the curvature runs along the non-radiating width; the
sagitta at the patch edge is $1.8$\,mm. A $40$\,mm air margin separates the
board from the UPML. The target is $|S_{11}| \le -10$\,dB over the
$2.400$--$2.4835$\,GHz ISM band.

\paragraph{The coarse mesh.}
The optimization mesh L1 has $150{,}561$ edges, a median of $2.0$ element
layers through the substrate under the patch and a mean edge length of
$0.78$\,mm at the patch rim. On L1 at $p = 1$ the initial patch's
$|S_{11}|$ dip is at $2.3118$\,GHz; at $p = 2$ on L1 (\PL) it is at
$2.3566$\,GHz, $45$\,MHz higher. A direct LU solve at \PL{} ($804{,}460$
unknowns) takes $132$\,GiB in geode-fem, so refining further at $p = 2$ is
not practical with this solver.

\begin{table}[t]
  \centering
  \caption{Where the error sits: the initial patch's $|S_{11}|$ dip under
  refinement of L1, geode-fem at $p = 1$. ``Slab'' is the element size in
  the substrate under the patch, ``rim'' the mean edge length within
  $0.5$\,mm of the patch edge, ``layers'' the median number of elements
  through the substrate under the patch. Model-dependent values.}
  \label{tab:x1}
  \small
  \begin{tabular}{@{}l c c c r c r@{}}
    \toprule
    Mesh & Slab (mm) & Rim (mm) & Layers & DOFs & Dip (GHz) & vs L1 (MHz) \\
    \midrule
    L1 & $0.80$ & $0.78$ & $2.00$ & $150{,}561$ & $2.3118$ & $0$ \\
    slab $0.533$ & $0.533$ & $0.52$ & $2.54$ & $234{,}691$ & $2.3273$ & $+15.5$ \\
    slab $0.40$ & $0.40$ & $0.39$ & $3.44$ & $414{,}837$ & $2.3404$ & $+28.6$ \\
    slab $0.30$ & $0.30$ & $0.29$ & $4.19$ & $828{,}775$ & $2.3480$ & $+36.2$ \\
    rim $0.40$ & $0.80$ & $0.39$ & $1.99$ & $317{,}424$ & $2.3420$ & $+30.2$ \\
    rim $0.20$ & $0.80$ & $0.20$ & $1.99$ & $430{,}808$ & $2.3569$ & $+45.1$ \\
    slab $0.30$ + rim $0.10$ & $0.30$ & $0.10$ & $4.19$ & $1{,}738{,}557$ & $2.3659$ & $+54.1$ \\
    \bottomrule
  \end{tabular}
\end{table}

\paragraph{The error sits at the rim.}
Table~\ref{tab:x1} refines L1 in two ways. Refining the whole slab under the
patch to $0.40$\,mm gives $3.4$ layers and, as a side effect, a $0.39$\,mm
rim. Refining only within $0.5$\,mm of the patch edge to the same rim size,
with the 2-layer interior unchanged, lands within $1.6$\,MHz of it
($2.3404$ against $2.3420$\,GHz). The interior layer count does not decide
the dip; the rim does, as expected where the field is singular at a
conductor edge \citep{meixner1972behavior}. A $0.2$\,mm rim at $p = 1$
($431$k unknowns) matches \PL{} ($805$k) to $0.3$\,MHz. Over all seven
$p = 1$ rungs the dip follows $f = f_\infty - C\,h_{\mathrm{rim}}^{\,0.80}$
with a $1.3$\,MHz rms residual (Figure~\ref{fig:accuracy}b). The exponent is
a fit to the rim edge length on graded meshes, not a demonstrated
asymptotic order. The combined rung, which halves the rim to $0.1$\,mm and
refines the slab, moves the dip $9.0$\,MHz.

\begin{figure}[t]
  \centering
  \anvilfig{fig2_accuracy.pdf}{fig2_accuracy.py}
  \caption{The initial patch's $|S_{11}|$ dip under refinement. (a) Against
  the number of unknowns: geode-fem at $p = 1$ on L1 and its slab, rim and
  combined refinements, geode-fem at $p = 2$ on L1, and Palace at
  $p = 1, 2, 3$ on the same L1 mesh. (b) Against the mean edge length at the
  patch rim, with the fit $f_\infty - C h^{0.80}$ (dashed), its limit and the
  order-1 Richardson step on the $0.39$/$0.20$\,mm rim pair. The shaded band,
  $2.366$--$2.380$\,GHz, runs from the finest computed value to the fit's
  limit; it is an estimate, not a bound. Data:
  \texttt{v6/v6\_summary.toml} \texttt{[x1]}.}
  \label{fig:accuracy}
\end{figure}

\paragraph{An independent solver on the same meshes.}
Palace \citep{palace} solved the same problem on geode-fem's own meshes,
written out with their physical tags, with the same materials, PEC surfaces
and uniform $50\,\Omega$ lumped port. Its absorber is different: Palace has
no UPML here, so the UPML shell is air and the outer wall carries a
first-order absorbing boundary. On L1, Palace and geode-fem agree on the
dip to $1.1$\,MHz at $p = 1$ ($2.3130$ against $2.3118$\,GHz) and $0.4$\,MHz
at $p = 2$ ($2.3570$ against $2.3566$\,GHz), and on the strip-fed design of
Section~\ref{sec:procedure} the \PL{} $-10$\,dB band reproduces to about
$1$\,MHz. Palace at $p = 3$ on L1 gives $2.3660$\,GHz, $8.9$\,MHz above its
$p = 2$ value and $0.02$\,MHz from geode-fem's finest rung (slab $0.30$ plus
rim $0.10$\,mm), $2.3659$\,GHz (Figure~\ref{fig:accuracy}a). Two independent
refinements of the same model, order in Palace and rim size in geode-fem,
meet at the same number.

\paragraph{The mesh-limit resonance.}
Every refinement computed so far raises the dip. If that continues, the
finest computed value, $2.366$\,GHz, is a lower estimate of the mesh-limit
dip; this is an inference from the trend, not a proof. The rim fit puts the
limit at $2.380$\,GHz and the order-1 Richardson step on the rim pair at
$2.372$\,GHz. We therefore estimate the initial patch's mesh-limit dip at
$2.366$--$2.380$\,GHz, $0$ to $14$\,MHz above the finest computed models.
The Hammerstad transmission-line estimate of the unloaded patch's resonance
is $2.433$\,GHz \citep{hammerstad1975equations}, about $50$\,MHz higher. The
two quantities differ: the probe's reactance pulls the matched dip below the
unloaded cavity resonance, and the closed form ignores the probe and the
finite ground. A port-free eigenfrequency would separate them and was not
computed.

\paragraph{Absorber and air margin.}
The UPML conductivity $\sigma_0 = 1$ was chosen from a ladder of $\sigma_0$
and shell thickness, over which $|S_{11}|$ changes by at most $0.021$ in
magnitude. Re-meshing the strip-fed design of Section~\ref{sec:procedure}
with a $60$ or $80$\,mm air margin instead of $40$\,mm moves its $p = 1$ dip
by at most $2.4$\,MHz and its $-10$\,dB band by at most $1.4$\,MHz, not
monotonically. The Palace runs, with a different absorber, agree with
geode-fem to the levels above. The absorber does not load the antenna at
the level of the refinement shifts.

\paragraph{Curvature.}
The same copper on a flat board (bend radius $10^6$\,mm, the same mesh node
for node) resonates $0.3$\,MHz from the bent one at $p = 1$, and a complex
frequency-shift fit at \PL{} finds $0.28$\,MHz. The bend does not detune the
patch. It changes its input impedance: on the strip-fed design the bent
board has a $93.3$\,MHz $p = 1$ $-10$\,dB band against $82.8$\,MHz flat, and
a worst in-band $|S_{11}|$ at \PL{} of $-10.78$ against $-9.93$\,dB. The
design was optimized on the bent board; a design optimized flat might
recover the difference, and was not run.

\section{The Gradient Against Gradient-Free Search}
\label{sec:baselines}

\paragraph{Setup.}
All optimizers solve one problem: the 38 parameters, the objective
(\ref{eq:objective}) on the unshifted band, the $0.25$ volume guard and the
L1 fixture at $p = 1$, starting from $\mathbf{c} = 0$ ($-1.47$\,dB). The
adjoint runs L-BFGS. Separable CMA-ES runs with population $\lambda = 14$
and initial step $\sigma = 0.5$ for 190 evaluations (one seed for 114), with
three seeds, and with geode-fem or Palace as the forward solver. NEWUOA runs
with Palace, initial trust radius $0.5$ (equal to $\sigma$), final radius
$10^{-3}$, PRIMA's default $2n + 1 = 77$ interpolation points and 190
evaluations. For Palace, geode-fem applies the morph and the guard, writes
the moved mesh with its tags, and Palace solves the five band frequencies
directly at $p = 1$ with GMRES on 16 MPI ranks. NEWUOA is unconstrained, so a
design that fails the guard receives a penalty above any feasible
objective without a solve; this happened once in 190 evaluations.

\paragraph{Counting.}
One solve is one band frequency with one matrix factorization. The adjoint
reuses the forward factorization for its adjoint solve and counts one solve
per frequency; its start, which yields its first gradient, is counted. The
gradient-free runs' separate evaluation at $\mathbf{c} = 0$ is not counted,
except as NEWUOA's own first evaluation. Guard-rejected CMA-ES samples cost
no solve. Each level is the worst of the five band points in that run's own
forward model. Wall-clock is on one host type, an r6i.4xlarge (16 vCPU):
measured for the Palace runs, which include mesh writes and Palace start-up,
and estimated for the geode-fem runs as solves times the per-solve cost
measured on that host, $14.4$\,s forward and $19.0$\,s forward plus adjoint.
A Palace band evaluation costs $64.4$\,s against geode-fem's $72.1$\,s, so
Palace's forward solve is about $11\%$ faster.

\begin{table}[t]
  \centering
  \caption{Solves to fixed worst-of-band levels on the L1 optimization model
  ($p = 1$), from $\mathbf{c} = 0$. Three CMA-ES entries are seeds (the
  first ran 114 evaluations, the others 190). ``Best'' is the best level in
  the run, ``in 125'' the best within the adjoint's 125 solves. Wall-clock
  to $-9.5$\,dB on an r6i.4xlarge: measured for Palace, estimated for
  geode-fem. A dash means the level was not reached.}
  \label{tab:baselines}
  \small
  \setlength{\tabcolsep}{4pt}
  \begin{tabular}{@{}l l c c c c c c@{}}
    \toprule
    Optimizer & Forward & \multicolumn{3}{c}{Solves to} & Best & In 125 & Wall to \\
    \cmidrule(lr){3-5}
    & & $-9$\,dB & $-9.5$\,dB & $-10$\,dB & (dB) & (dB) & $-9.5$\,dB (min) \\
    \midrule
    L-BFGS, adjoint & geode-fem & $30$ & $35$ & $60$ & $-10.04$ & $-10.04$ & $11.1$ (est.) \\
    sep-CMA-ES & geode-fem & $405$ & $520$ & -- & $-9.56$ & $-2.66$ & $125$ (est.) \\
    & & $385$ & $385$ & $695$ & $-10.02$ & $-2.40$ & $93$ (est.) \\
    & & $395$ & $395$ & -- & $-9.77$ & $-2.82$ & $95$ (est.) \\
    sep-CMA-ES & Palace & $405$ & $805$ & -- & $-9.58$ & $-2.66$ & $172$ \\
    & & $385$ & $385$ & -- & $-9.89$ & $-2.40$ & $84$ \\
    & & $395$ & $570$ & -- & $-9.90$ & $-2.83$ & $124$ \\
    NEWUOA (PRIMA) & Palace & $415$ & $420$ & -- & $-9.95$ & $-1.52$ & $90$ \\
    \bottomrule
  \end{tabular}
\end{table}

\begin{figure}[t]
  \centering
  \anvilfig{fig3_optimizers.pdf}{fig3_optimizers.py}
  \caption{Best worst-of-band $|S_{11}|$ so far on the L1 optimization model,
  for the adjoint L-BFGS and the gradient-free runs of
  Table~\ref{tab:baselines}, (a) against solves and (b) against wall-clock
  on one r6i.4xlarge. Palace wall-clock is measured; geode-fem wall-clock is
  estimated from the measured per-solve cost on that host. Data:
  \texttt{copper/L1\_copper\_38.toml}, \texttt{v6/x4/}, \texttt{v6/palace\_loop/}.}
  \label{fig:optimizers}
\end{figure}

\paragraph{Results.}
Table~\ref{tab:baselines} and Figure~\ref{fig:optimizers} give the
comparison. To reach $-9.5$\,dB the adjoint uses $35$ solves; CMA-ES uses
$385$--$520$ with geode-fem and $385$--$805$ with Palace, and NEWUOA $420$,
that is $11$ to $23$ times as many. In wall-clock the Palace-driven runs take
$84$--$172$\,min (CMA-ES) and $90$\,min (NEWUOA) against an estimated
$11.1$\,min for the adjoint, $7.6$ to $15.5$ times as long. The wall-clock
ratio is smaller than the solve ratio because an adjoint evaluation costs
$1.32$ forward evaluations and Palace's forward solve is the faster one.
Within the adjoint's whole budget of 125 solves, the gradient-free runs are
between $-1.52$ and $-2.83$\,dB. NEWUOA spends its first 77 evaluations
(385 solves) building its interpolation model, which with $2n + 1$ points is
close to one finite-difference gradient, and then steps to $-9.61$\,dB within
8 evaluations; a smaller interpolation set was not tried.

The adjoint reaches $-10$\,dB in $60$ solves. One geode-fem CMA-ES seed of
three reaches it, at $695$, and then matches the adjoint's final value
($-10.02$ against $-10.04$\,dB); no Palace-driven run reaches it within 190
evaluations ($660$ to $945$ solves). The $-10$\,dB column needs care. The two
forward models agree to $0.004$\,dB on the 41 designs both evaluated (the
first CMA-ES generation of each seed, 42 samples with one rejected by the
guard in both), but the difference grows at optimized designs, where the
band edge sits on a steep slope: Palace reads the adjoint's final design at
$-9.80$\,dB against geode-fem's $-10.04$\,dB, and geode-fem reads the four
Palace-optimized designs at $-9.81$ to $-9.89$\,dB, within $0.25$\,dB of
Palace's own values. A $0.2$\,dB model difference is the size of the gap
between reaching and missing $-10$\,dB, so that column compares two slightly
different targets. The $-9$ and $-9.5$\,dB columns do not depend on it.

\paragraph{What the gradient buys, and what it does not.}
This is an ablation of the gradient: the geode-fem CMA-ES runs share the
adjoint's forward model, objective, basis and guard, and the Palace runs
show that a faster forward solver does not close the gap. On this problem
the gradient buys speed, about an order of magnitude in solves, and not a
better optimum. Each gradient-free optimizer was run at one setting
($\sigma = 0.5$, $\lambda = 14$; PRIMA's defaults), so a tuned gradient-free
method would do better than these rows. All levels are on the L1 model at
$p = 1$, which sits about $45$--$57$\,MHz below the refined models
(Section~\ref{sec:accuracy}); none is a level of the refined model.

\paragraph{A four-parameter rectangle.}
The 38-parameter basis is not needed to meet the target. Collapsing it to
the patch width and length plus the two feed parameters, with the same
objective, guard and band shift as the strip-fed design's final step and
the same feed re-basing (Section~\ref{sec:procedure}), the adjoint widens
the patch from $38$ to $46.1$\,mm and reaches a worst in-band $|S_{11}|$ of
$-10.33$\,dB at \PL{}, against $-10.78$\,dB for the 38-parameter design on
the same model. Its $p = 1$ $-10$\,dB band is $86.9$\,MHz against
$93.3$\,MHz, so the free outline adds about $6$\,MHz of bandwidth. The
rectangle started with the 38-parameter chain's final band shift instead of
finding it, so its 340 driven solves are not comparable with the chain's
655, which include the steps that found the shift.

\section{Design Procedure}
\label{sec:procedure}

\paragraph{Frequency space mapping.}
A design optimized on L1 resonates higher on a finer model. We model the
fine response as the coarse response shifted in frequency,
$S_{\mathrm{fine}}(f) \approx S_{\mathrm{coarse}}(f - \delta)$, fit $\delta$
by least squares over the fine model's points, and re-optimize on the coarse
model for the band shifted down by $\delta$. For the cold L1 optimum of
Section~\ref{sec:baselines} the fitted L1-to-\PL{} shift is $47.85$\,MHz, and
the root-mean-square complex misfit falls from $0.23$--$0.38$ unshifted to
$0.018$--$0.025$ shifted over the designs fitted. That shift is about
$3$\,MHz from the initial patch's own L1-to-\PL{} shift ($44.8$\,MHz), so the
copper motion adds little detune of its own. Without the correction, the
cold L1 optimum ($-10.04$\,dB on L1) gives $-4.76$\,dB at \PL. This
ablation is confounded, because the corrected chain also re-bases the feed
and widens the patch; a cleaner comparison is a design optimized on the
finer mesh L2 with a $30$\,MHz shift and no re-base, which gives
$-9.92$\,dB at \PL, against a fitted L2 shift of $29.33$\,MHz.

\paragraph{Feed re-basing.}
Moving the feed toward the $+y$ edge compresses the mesh between them, and
the cold runs stop there at the volume guard. We then regenerate the mesh
with the port at its moved position and carry the design over by shifting
the two feed coefficients. After re-basing, the optimizer widens the patch
until the guard binds in the air $5.5$\,mm outside a $\pm s$ edge, which is
a non-radiating edge. The widening, not the feed, is where every later
descent stops.

\paragraph{Three designs.}
Table~\ref{tab:chain} lists the designs. Design~A is the end of the
space-mapping chain centred for \PL{} with the feed width free, which widened
the port into a $4.27$\,mm strip. At \PL{} it reaches $-10.78$\,dB with a
$94.9$\,MHz band, $2.3938$--$2.4887$\,GHz. At Palace $p = 3$ the band moves
up to $2.4029$--$2.4975$\,GHz, the lower edge misses $2.400$\,GHz by
$2.9$\,MHz, and the worst in-band value is $-9.56$\,dB: a design centred on
\PL{} misses once the rim is resolved. Design~B re-runs A's last step with
the probe width pinned at $1.6$\,mm and reaches $-10.62$\,dB at \PL{}
($-10.65$\,dB with a $1.27$\,mm probe), so the wide strip was not what
carried the bandwidth. Design~C re-centres B for the refined model with a
band shift of $57.0$\,MHz: $44.94$\,MHz from L1 to \PL, $8.7$\,MHz for design~A's
$p = 2$ to $p = 3$ dip shift in Palace, and
$3.4$\,MHz further down because every refinement so far has moved the band
up. A parallel run at $53.6$\,MHz was set aside: adding the initial patch's
step of about $9$\,MHz from the $0.2$\,mm-rim mesh to Palace $p = 3$ to each
design's lower edge on that mesh predicts a lower-edge margin at $p = 3$ of
about $2.6$\,MHz for it and about $6$\,MHz for C. The descent to C ran 10 steps and 125 solves,
from $-8.73$ to $-10.64$\,dB on its shifted L1 band, and stopped on the guard
with the gradient norm down from $0.143$ to $2.2\times10^{-3}$; the gradient
matched finite differences to $1.2\times10^{-6}$. As built, C is
$29.00$\,mm long and $49.72$\,mm wide on the centre line, with a barrel
outline whose $\pm s$ edges bow out most at mid-length, and its $1.6$\,mm
probe sits $5.14$\,mm from the $+y$ edge.

\begin{table}[t]
  \centering
  \caption{The designs. $\delta$ is the band shift of the final descent.
  ``Solves'' counts the driven solves of that descent as recorded,
  including its gradient checks. ``Fine model'' is the model the design was
  centred for.}
  \label{tab:chain}
  \small
  \begin{tabular}{@{}l l c c c l@{}}
    \toprule
    Design & Feed & $\delta$ (MHz) & Steps & Solves & Fine model \\
    \midrule
    A & $4.27$\,mm strip (free width) & $44.94$ & $9$ & $125$ & \PL \\
    B & $1.6$\,mm probe (pinned) & $44.94$ & $10$ & $130$ & \PL \\
    C (final) & $1.6$\,mm probe (pinned) & $57.0$ & $10$ & $125$ & refined rim, Palace $p = 3$ \\
    \bottomrule
  \end{tabular}
\end{table}

\section{Results}
\label{sec:results}

\begin{figure}[t]
  \centering
  \anvilfig{fig4_response.pdf}{fig4_response.py}
  \caption{(a) Return loss of design~C. Palace at $p = 3$ and $p = 2$ on the
  design's L1 mesh (1\,MHz grids); geode-fem at $p = 1$ on the combined
  $0.1$\,mm-rim and $0.3$\,mm-slab mesh (13 points) and on the $0.2$\,mm-rim
  mesh; and geode-fem at $p = 1$ on L1, the optimization model, on which the
  design sits $57$\,MHz low by construction. The dashed line is $-10$\,dB and
  the shaded band is the ISM band. (b) $|S_{11}|$ at $2.400$\,GHz if the
  whole Palace $p = 3$ response moves up by a further $d$, read from the
  measured response at $2.400 - d$\,GHz, with the two drift estimates of
  Section~\ref{sec:accuracy} marked. Data: \texttt{recentre/}.}
  \label{fig:response}
\end{figure}

\begin{table}[t]
  \centering
  \caption{Design~C on each model, in order of refinement, with design~A on
  Palace $p = 3$ for comparison. ``Worst'' is over the nine in-band points
  ($2.400$ to $2.4835$\,GHz in $10.4375$\,MHz steps). Band crossings are
  interpolated linearly in dB. Margins are $2.400$\,GHz minus the lower
  crossing and the upper crossing minus $2.4835$\,GHz. Model-dependent
  values.}
  \label{tab:verify}
  \small
  \setlength{\tabcolsep}{4pt}
  \begin{tabular}{@{}l r c c c c@{}}
    \toprule
    Model & DOFs & Worst (dB) & $-10$\,dB band (GHz) & Coverage & Margins (MHz) \\
    \midrule
    geode-fem $p = 1$, L1 & $150{,}541$ & $-4.94$ & $2.3393$--$2.4304$ & $36\%$ & $60.7$ / $-53.1$ \\
    geode-fem $p = 2$, L1 & $804{,}336$ & $-9.08$ & not swept & not swept & not swept \\
    Palace $p = 2$, L1 & $884{,}536$ & $-9.08$ & $2.3847$--$2.4770$ & $92.2\%$ & $15.3$ / $-6.5$ \\
    geode-fem $p = 1$, rim $0.2$ & $431{,}528$ & $-9.23$ & $2.3849$--$2.4782$ & $93.6\%$ & $15.1$ / $-5.3$ \\
    geode-fem $p = 1$, rim $0.1$ + slab & $1{,}738{,}279$ & $-10.73$ & $2.3934$--$2.4881$ & $100\%$ & $6.6$ / $4.6$ \\
    Palace $p = 3$, L1 & $2{,}503{,}077$ & $-10.46$ & $2.3935$--$2.4864$ & $100\%$ & $6.5$ / $2.9$ \\
    \midrule
    Design A, Palace $p = 3$ & $2{,}729{,}331$ & $-9.56$ & $2.4029$--$2.4975$ & $96.5\%$ & $-2.9$ / $14.0$ \\
    \bottomrule
  \end{tabular}
\end{table}

\paragraph{On the finest models.}
On Palace at $p = 3$ design~C has a worst in-band $|S_{11}|$ of
$-10.46$\,dB, at $2.4835$\,GHz, and a contiguous $-10$\,dB band from
$2.3935$ to $2.4864$\,GHz, $93.0$\,MHz wide, which covers the whole
$83.5$\,MHz ISM band with $6.5$\,MHz to spare at $2.400$\,GHz and $2.9$\,MHz
at $2.4835$\,GHz (Table~\ref{tab:verify}, Figure~\ref{fig:response}a). At
$2.400$\,GHz it is at $-11.16$\,dB. geode-fem's finest mesh, the combined
$0.1$\,mm rim and $0.3$\,mm slab ($1.74$M unknowns), gives $-10.73$\,dB and a
$94.7$\,MHz band, $2.3934$--$2.4881$\,GHz. The two lower edges agree to
$0.02$\,MHz, as the two codes agreed on the initial patch's dip. The upper
edges differ by $1.7$\,MHz, partly because geode-fem's upper crossing is
interpolated across a $5$\,MHz gap and Palace's across $1$\,MHz. The two
models are not equally independent of the design. Palace $p = 3$ informed
the band shift, through design~A's $p = 2$ to $p = 3$ step; geode-fem's
combined mesh was not used in choosing the design. On the coarser models the
design sits high by construction and misses at the upper edge: at \PL{}
Palace and geode-fem agree to $0.003$\,dB on its worst in-band value
($-9.076$ and $-9.079$\,dB). From $p = 2$ to $p = 3$ the design's dip rises
$9.1$\,MHz and its lower band edge $8.7$\,MHz; the initial patch's dip rises
$8.9$\,MHz over the same step, so the design shifts like the patch.

The refined geode-fem meshes were made with a rim refinement laid around
the initial $38\times29$\,mm outline; measured on the moved meshes, the mean
rim edge length around the moved outline is $0.21$\,mm and $0.10$\,mm,
unchanged by the design. On the combined mesh the default air morph
inverts elements $0.1$\,mm above the patch edge, so the air there was moved
on cylinders instead, which keeps the copper, slab and ground identical; on
the $0.2$\,mm-rim mesh the two air motions agree to $0.09$\,MHz at the lower
edge and $0.015$\,dB in the worst in-band value.

\paragraph{Margin against the remaining drift.}
If the whole band moves up by a further $d$, design~C keeps the ISM band
covered at the lower edge while $d \le 6.5$\,MHz; the upper edge only gains.
The estimates of the remaining drift above the finest models are those of
the initial patch (Section~\ref{sec:accuracy}). For the order-1 Richardson
estimate, $d = 6.4$\,MHz, the lower edge reaches about $2.3999$\,GHz and
$2.400$\,GHz sits at $-10.0$\,dB, just inside. For the rim-fit estimate,
$d = 13.5$\,MHz, the lower edge reaches about $2.4070$\,GHz and
$2.400$\,GHz sits at $-8.94$\,dB, a miss (Figure~\ref{fig:response}b).
These two values are the measured $p = 3$ response shifted, not computed.
So the design meets the target with margin on the most accurate computed
models, stays inside the band for the lower part of the drift estimate, and
misses at $2.400$\,GHz if the drift reaches the rim-fit value. Covering the
whole $0$ to $14$\,MHz range would need about $97.5$\,MHz of $-10$\,dB
bandwidth, $4.5$\,MHz more than design~C has.

\paragraph{The guard limits the bandwidth.}
The descents that produced designs A, B and C each stop on the $0.25$
volume guard in the air outside the widening patch, not at a stationary point of the objective, so
the $91$--$95$\,MHz bandwidth and with it the slack are what the guard
allows. Regenerating the mesh at the widened outline would relieve the
guard; it was not run.

\paragraph{Radiation efficiency and gain.}
For the strip-fed design~A, whose outline matches C's to $0.2$\,mm in length
and $0.03$\,mm in width, we computed the accepted port power, the Poynting
flux through a closed surface in the air box, the FR-4 loss and the
broadside gain by a near-to-far-field transform, at $p = 1$ on L1 and on the
finer L3 mesh. The power balance closes to $2.3\%$ on L1 and $1.2\%$ on L3.
The radiation efficiency is $0.61$--$0.62$ on L1 and $0.62$ on L3 across the
band, the dielectric loss fraction $0.36$, and the broadside gain
$4.4$--$4.6$\,dBi with a front-to-back ratio of $12$--$15$\,dB. With PEC
conductors this efficiency counts radiation against dielectric loss only;
copper loss would lower it. These values were not recomputed for
design~C, whose feed differs.

\paragraph{Bandwidth against a closed form.}
On the voltage-standing-wave-ratio-2 criterion, Palace gives design~A a
bandwidth of $100.3$\,MHz at $p = 2$ and $101.0$\,MHz at $p = 3$ (design~C:
$99.0$\,MHz at $p = 3$). The closed-form formulas of
\citet{jackson1991simple}, in the form of Jackson's course notes
\citep{jackson_ece6345}, give $70.4$\,MHz for a $49.75\times28.80$\,mm
rectangle on this substrate, and $67.9$\,MHz at the area-equivalent width of
$45.4$\,mm. The gap is in the radiation $Q$. The only loss in the model is
$\tan\delta$, so the FEM's dielectric loss fraction bounds its total $Q$
from below, $Q_t \ge 0.362/\tan\delta = 18.1$, and with $\eta = 0.615$ its
radiation $Q$ is at least $29$; the closed form gives $Q_t = 24.5$ and a
radiation $Q$ of $48.0$, with radiated and dielectric fractions of
$0.51$ and $0.49$ against the FEM's $0.62$ and $0.36$. The model patch
radiates more strongly than the infinite-ground closed form; the finite
$62\times53$\,mm ground is a plausible cause that we did not test.

\paragraph{Cost.}
The final descent took $26$\,min on a 32-vCPU cloud host shared with a
second run. The Palace $p = 3$ sweep of design~C (281 frequencies from an
adaptive reduced-order model built on 12 full solves at $2.5$M unknowns) took
$3.5$\,h on 16 MPI ranks. On the combined mesh geode-fem's direct solver
took about $8$\,min per frequency at $204$\,GiB peak memory, which is why
that mesh served for verification only.

\section{Limitations and Future Work}
\label{sec:limits}

\paragraph{Feed and losses.}
The probe is a uniform lumped port of the pin's width across the substrate.
A coaxial feed with its own reactance is not modelled, conductors are PEC,
and the design has not been fabricated or measured.

\paragraph{Refinement.}
The remaining drift is estimated, not computed, and design~C covers only
the lower part of its range. The adjoint is first order; a $p = 2$ shape
adjoint, or optimizing on a rim-refined mesh directly, would remove the
coarse-to-fine shift that space mapping corrects here. A port-free
eigenfrequency, which would separate the feed reactance from the patch
resonance, was not computed, and Palace ran on L1 only.

\paragraph{Scope of the comparison.}
The optimizer comparison is one problem, one starting point, one setting
per gradient-free optimizer and one forward model level. A smaller NEWUOA
interpolation set, Bayesian optimization or a tuned CMA-ES might narrow the
gap. The comparison measures search cost on the coarse model; it does not
test whether the gradient helps at the refined level.

\paragraph{Comparison with FDTD.}
A matched-geometry Meep \citep{oskooi2010meep} model of the initial antenna,
run at three grid resolutions down to $0.25$\,mm cells, disagrees with
geode-fem and Palace. In the sampled $2.2$--$2.7$\,GHz window its only deep
$|S_{11}|$ feature sits at the $2.20$\,GHz lower edge, so its resonance is at
or below $2.20$\,GHz, while geode-fem and Palace place the dip at or above
$2.357$\,GHz; and under refinement the two move in opposite directions, the
finite-element dip upward and Meep's minimum downward. The cause has not
been identified and is being diagnosed. Until it is, we make no comparison of
cost or accuracy with FDTD, including a conformal FDTD scheme
\citep{dey1997locally,zagorodnov2003uniformly}.

\section{How geode-fem Was Designed and Built}
\label{sec:built}

geode-fem implements the whiteroom operator specification of Crutcher
Dunnavant \citep{dunnavant2026whiteroom}, whose four layers restate Palace's
cited source as pure functions, fused primitives, tensor operations and a
small graph-evaluation calculus. He also gave design input on the solver's
operator layer and contributed code: the migration to Rust 2024, the Metal
(Apple GPU) backend, the present layout of the core crate and the examples,
and the geode-util crate. The rest of the code, the experiments and the text
of this paper were produced by AI models working as role-separated agents:
Anthropic's Claude (Opus, Fable and Sonnet models, through Claude Code) and
OpenAI's GPT Sol and GPT Astra models (via Codex). The agents were
coordinated by Loom \citep{loom}, an issue, worktree and pull-request
workflow, and by Squad, a message room the agents share. Nearly every
change was reviewed by a separate agent session. The agent reviews of the
experiments behind this paper recomputed the reported numbers from the
archived records, and the script that checks this paper's text against
those records was mutation-tested. The human author directed the research,
made every scoping and claim decision, and takes responsibility for the
content.

\section{Reproducibility}
\label{sec:repro}

Every result in this paper is read from an archived, scripted
artifact.\footnote{The geode-fem source is public at
\url{https://github.com/rjwalters/geode-fem}. Archival DOI to be added at
deposit.}
% TODO(operator): Zenodo DOI
Table~\ref{tab:artifacts} lists them, relative to
\path{benchmarks/patch_antenna_conformal/}. Each run record holds the exact
command and the SHA-256 of its mesh; the meshes are regenerated from YAML
recipes in \path{reference/gmsh/} with gmsh 4.15.2, and runs on re-based
fixtures need the \texttt{--feed-support-legacy} flag. Most records carry no
generating commit and one carries an intermediate one; the code that
produced them is merged at geode-fem commit \texttt{99646f06}. Palace is
commit \texttt{fba6a5b}, built from
\path{reference/palace/docker/Dockerfile}. PRIMA was built from source at
libprima commit \texttt{43863c69} (Python package
\texttt{0.7.3.dev1801+g43863c69c}); the compiler was not recorded. The
geode-fem solves ran on AWS Graviton4 hosts (r8g.4xlarge and r8g.8xlarge),
and the Palace solves and the cost measurements on r6i.4xlarge. The figure
scripts in \path{figures/src/} read the artifacts directly,
\path{provenance.md} in this paper's directory maps every number to its
source field, and \path{verify_evidence.py} checks the text against them.

\begin{table}[t]
  \centering
  \caption{Artifacts behind each result, relative to
  \texttt{benchmarks/patch\_antenna\_conformal/}.}
  \label{tab:artifacts}
  \footnotesize
  \begin{tabular}{@{}>{\raggedright\arraybackslash}p{0.30\linewidth}>{\raggedright\arraybackslash}p{0.64\linewidth}@{}}
    \toprule
    Result & Artifact \\
    \midrule
    Fixture, coarse ladder, absorber ladder & \path{baseline_2g4_L1.toml}, \path{convergence.toml}, \path{upml_ladder_L1.toml} \\
    Rim and slab refinement, Palace cross-check, curvature, radiation, air margin (Section~\ref{sec:accuracy}) & \path{v6/v6_summary.toml} from \path{scripts/v6_summary.py}; \path{v6/x1/}, \path{v6/x2/}, \path{v6/x3/} \\
    Gradient-free baselines (Section~\ref{sec:baselines}) & \path{v6/palace_loop/palace_loop_summary.toml} from \path{scripts/palace_loop_summary.py}; \path{v6/x4/}, \path{v6/palace_loop/}; \path{copper/L1_copper_38.toml} \\
    Space mapping, designs A and B (Section~\ref{sec:procedure}) & \path{copper/}, \path{copper/g2_summary.toml}, \path{scripts/copper_shift_fit.py}; \path{v6/x5/} \\
    Design C and its checks (Section~\ref{sec:results}) & \path{recentre/recentre_summary.toml} from \path{scripts/recentre_summary.py}; \path{recentre/}, \path{recentre/palace/} \\
    FDTD comparison (Section~\ref{sec:limits}) & \path{reference/meep/docker/AWS_R4.md}; \path{benchmarks/fdtd_density_baseline/meep_runs/aws_r4/resolution_trend.json} \\
    Generator and tests & \path{crates/geode-core/examples/patch_conformal_nested_dof.rs}; \path{crates/geode-core/tests/patch_conformal_*.rs} \\
    \bottomrule
  \end{tabular}
\end{table}

\section{Conclusion}
\label{sec:conclusion}

On a 38-parameter copper design of a bent 2.4\,GHz patch, an adjoint
gradient reaches the band target with an order of magnitude fewer solves
than two gradient-free optimizers, whether they are driven by the same
solver or by a faster independent one. The gradient buys speed, not a
better optimum, and a four-parameter rectangle would also have met the
target. Getting the target to hold on a fine model needed the forward model
resolved at the patch rim, where an independent solver confirms where the
resonance sits, and one more space-mapping step. The resulting design,
with a $1.6$\,mm probe on a standard board, covers the ISM band at
$-10$\,dB on both fine models with $6.5$\,MHz of margin at the lower edge.
That margin covers the order-1 estimate of the remaining drift but not the
rim-fit estimate, and the volume guard, not the objective, limits it.

\paragraph{Acknowledgments.}
We thank Crutcher Dunnavant for the whiteroom specification, for his
design input to the solver's operator layer and for his contributions to
the codebase. The AI-orchestrated development of geode-fem and of this
paper is described in Section~\ref{sec:built}.
% TODO(operator): approve acknowledgment wording

\bibliographystyle{plainnat}
\bibliography{refs}

\end{document}
